\scanpage{154}
\[
 \boxed{\begin{aligned}
 (x_1,x_2)\longmapsto x_1\ast x_2&=t_1\otimes x_2+x_1\otimes t_2-2x_1\otimes x_2\\
 &=\sigma_1(x_1)\otimes x_2+x_1\otimes\sigma(x_2)
 \end{aligned}}
\]
% slanted
\margin{\uline{NB} Il y a lieu de le prolonger de façon bil. aux modules enveloppants $E_1,E_2$ par $x_1\ast x_2=\underset{\varepsilon(x_1)t_1-x_1}{\underbrace{\sigma_1(x_1)}}\otimes x_2+\deleted{\unsure{\sigma}}x_1\otimes\underset{\varepsilon_2(x_2)t_2-x_2}{\underbrace{\sigma_2(x_2)}}=\varepsilon(x_1)t_1\otimes x_2+\varepsilon(x_2)\unsure{t_2\otimes x_2|x_1\otimes t_2}-2x_1\otimes x_2$.}
on vérifie que sur $A_1^0\times A_2^0$, cette application $f_{A_1^{\unsure{0}},A_2}$, avec
\[
 \begin{aligned}
 f^0_{A_1,A_2}: A_1^0\times A_2^0&\longrightarrow A_1^0\otimes A_2^0\\
 (\alpha_1,\alpha_2)&\longmapsto\alpha_1\otimes\alpha_2
 \end{aligned}
\]

\deleted{\unsure{Ceci}} Bien sûr, $A_1\ast A_2$ est un bifoncteur en $A_1,A_2$. Donc il y a deux involutions $\sigma_1\ast\mathrm{id}_{A_2}$ et $\mathrm{id}_{A_1}\ast\sigma_2$ sur $A_1\ast A_2$, définies par
\[
 x_1\ast x_2\longmapsto\sigma_1(x_1)\ast x_2\qquad
 x_1\ast x_2\longmapsto x_1\ast\sigma_2(\unsure{x_2|\alpha_2})
\]
Je dis qu'elles sont égales, et qu'on a
\[
 \underbrace{\sigma_1(x_1)}_{t_1-x_1}\ast x_2=x_1\ast\underbrace{\sigma_2(x_2)}_{t_2-x_2}=t_1\otimes t_2-x_1\ast x_2
\]
% slanted
\margin{Si $2\,1_k=0$, on a $t_1\otimes t_2\in A_1^0\otimes A_2^0$ ds espace des translations c'est celui qui définit l'inv. can. de $A_1\ast A_2$}
Ainsi, $A_1\ast A_2$ est muni aussi d'une involution \inserted{$\sigma_{A_1\ast A_2}$}. Il est solution du pb universel, cette fois dans les espaces affines involutifs sur $k$, pour les $\sigma$-accouplements d'espaces affines --- où cette fois on met la condition supplémentaire sur $f$
\[
 f(\sigma_1(x_1),x_2)=f(x_1,\sigma_2(x_2))=\sigma_B f(x_1,x_2).
\]
% archive sheet number: 153

